\documentclass[11pt]{article}
\usepackage[T1]{fontenc}
\usepackage[utf8]{inputenc}
% one typeface for everything: Palatino for text and mathematics, with real small caps
\usepackage[sc]{mathpazo}
\linespread{1.05}
\usepackage{amsmath,amssymb,amsthm}
\usepackage{microtype}
\usepackage[margin=1.1in]{geometry}
\usepackage[hidelinks,pdftitle={Two operations and a machine that looks for surprises},
  pdfsubject={An introduction from scratch to the duet project}]{hyperref}

\newcommand{\op}{\cdot}
\newcommand{\cop}{\circ}
\newcommand{\N}{\mathbb{N}}
\newcommand{\Z}{\mathbb{Z}}
\newcommand{\proj}{\textsc{duet}}
\newcommand{\file}[1]{\textit{#1}}          % file names: italic, same typeface

% headings: only two sizes (large for sections, normal for run-in paragraphs), all bold, same typeface
\makeatletter
\renewcommand\section{\@startsection{section}{1}{\z@}{-3.2ex \@plus -1ex \@minus -.2ex}{1.4ex \@plus .2ex}%
  {\normalfont\large\bfseries}}
\renewcommand\paragraph{\@startsection{paragraph}{4}{\z@}{1.8ex \@plus .5ex \@minus .2ex}{-0.6em}%
  {\normalfont\normalsize\bfseries}}
% a part is a quiet small-caps line, not a page-sized heading
\newcommand{\partline}[2]{\par\bigskip\addcontentsline{toc}{part}{Part #1\quad #2}%
  {\centering\normalfont\large\scshape Part #1\quad #2\par}\nobreak\medskip}
% table of contents: parts in small caps at normal size, sections without extra space between them
\renewcommand*\l@part[2]{\addpenalty{-\@highpenalty}\vskip .8em \@plus\p@
  {\parindent\z@\rightskip\@pnumwidth\parfillskip-\@pnumwidth\normalfont\scshape #1\hfil\hb@xt@\@pnumwidth{\hss #2}\par}}
\let\intro@lsection\l@section
\renewcommand{\l@section}[2]{\begingroup\let\addvspace\@gobble\intro@lsection{#1}{#2}\endgroup}
\makeatother

% every statement upright, with a bold label
\newtheoremstyle{uniform}{\medskipamount}{\medskipamount}{\normalfont}{}{\bfseries}{.}{0.5em}{}
\theoremstyle{uniform}
\newtheorem{theorem}{Theorem}[section]
\newtheorem{proposition}[theorem]{Proposition}
\newtheorem{definition}[theorem]{Definition}
\newtheorem{example}[theorem]{Example}
\newtheorem{exercise}{Exercise}

% Two versions from one source. intro.pdf: exercises at the end of each section. intro-reading.pdf (built with
% \def\readingversion{} first, see the Makefile): no exercises in the main text, which never refers to them;
% exercises and solutions in two appendices.
\newif\ifreading
\ifdefined\readingversion \readingtrue \else \readingfalse \fi
\newcommand{\either}[2]{\ifreading#2\else#1\fi}      % {text with exercises}{reading text}
\newcommand{\exnote}[1]{\ifreading\else#1\fi}        % a pointer to an exercise, only with exercises
\newcounter{exgroup}
\renewcommand{\theexercise}{\arabic{exgroup}.\arabic{exercise}}
\newcommand{\exgroup}[2]{\setcounter{exgroup}{#1}\setcounter{exercise}{0}%
  \ifreading\par\bigskip\noindent{\bfseries Section #1\quad #2}\par\nobreak\fi}
\newcommand{\inlineexercises}[1]{\ifreading\else\csname exercises@#1\endcsname\fi}

% an exercise points to its solution, the solution points back
\newcommand{\solref}[1]{\hfill\mbox{\textit{(solution on p.~\pageref{sol:#1})}}}
\newcommand{\solution}[1]{\par\medskip\noindent\textbf{Solution~\ref{ex:#1}.}\label{sol:#1}
  \textit{(exercise on p.~\pageref{ex:#1})}\ }

% Exercises of intro.tex, one macro per section. Inline at the end of each section in intro.pdf;
% collected in an appendix in intro-reading.pdf. Solutions are in intro.tex (appendix) in both versions.

% Section 1: Sets and functions
\expandafter\newcommand\csname exercises@1\endcsname{%
\exgroup{1}{Sets and functions}%
\begin{exercise}\label{ex:prod}
Write out $\{0,1\}\times\{0,1,2\}$. How many elements does it have? \solref{prod}
\end{exercise}
\begin{exercise}\label{ex:funcs}
How many functions are there from $\{0,1,2\}$ to $\{0,1\}$? How many are injective? How many are surjective?
\solref{funcs}
\end{exercise}
\begin{exercise}\label{ex:inj}
Check Proposition~\ref{prop:finite} on $A=\{0,1,2\}$ by counting: how many functions $A\to A$ are injective, and how
many of those are not surjective? \solref{inj}
\end{exercise}
}

% Section 2: Binary operations
\expandafter\newcommand\csname exercises@2\endcsname{%
\exgroup{2}{Binary operations}%
\begin{exercise}\label{ex:count}
Check the counts $16$ and $19{,}683$. \solref{count}
\end{exercise}
\begin{exercise}\label{ex:idem2}
An operation is \emph{idempotent} if $x\op x=x$ for every $x$. How many of the 16 operations on $\{0,1\}$ are
idempotent? Name them. \solref{idem2}
\end{exercise}
}

% Section 3: Laws
\expandafter\newcommand\csname exercises@3\endcsname{%
\exgroup{3}{Laws}%
\begin{exercise}\label{ex:max}
Which of these laws hold for $\max$ on $\{0,1,2\}$: associativity, commutativity, idempotence, $x\op y=x$,
$x\op(x\op y)=x\op y$? Is there an element $e$ with $e\op x=x$ for all $x$? \solref{max}
\end{exercise}
\begin{exercise}\label{ex:lp}
Show that the left projection $x\op y=x$ on a set with at least two elements is associative and idempotent but
not commutative. \solref{lp}
\end{exercise}
\begin{exercise}\label{ex:cost}
How many assignments must you try to check associativity on 3 elements? On 6? \solref{cost}
\end{exercise}
\begin{exercise}\label{ex:order1}
List the laws of order~1 in one operation, up to renaming. (One side is a variable, the other a product of two
variables.) \solref{order1}
\end{exercise}
}

% Section 4: Theories, models and implication
\expandafter\newcommand\csname exercises@4\endcsname{%
\exgroup{4}{Theories, models and implication}%
\begin{exercise}\label{ex:abs}
The \emph{absorption laws} for two operations are $x\op(x\cop y)=x$ and $x\cop(x\op y)=x$. Prove that together they
imply $x\op x=x$. (Hint: use the second law with $y=x$, then the first law with $y=x\op x$.) \solref{abs}
\end{exercise}
\begin{exercise}\label{ex:abs2}
How many pairs of operations $(\op,\cop)$ on $\{0,1\}$ satisfy both absorption laws? How many of them are
non-degenerate? Here that means both operations combine their arguments and the two operations differ.
\solref{abs2}
\end{exercise}
\begin{exercise}\label{ex:abs3}
Show that the first absorption law \emph{alone} does not imply $x\op x=x$: find a model of $x\op(x\cop y)=x$ on
$\{0,1\}$ in which $x\op x=x$ fails. \solref{abs3}
\end{exercise}
\begin{exercise}\label{ex:corpus}
Suppose you only know two structures: $\max$ on $\{0,1\}$ and the left projection on $\{0,1\}$. Both are idempotent
and both are associative. Does idempotence imply associativity? Test $x\op y=2x+2y \bmod 3$ on $\{0,1,2\}$.
\solref{corpus}
\end{exercise}
}

% Section 5: Two operations
\expandafter\newcommand\csname exercises@5\endcsname{%
\exgroup{5}{Two operations}%
\begin{exercise}\label{ex:lat3}
How many pairs $(\op,\cop)$ on $\{0,1,2\}$ are lattices, counting tables, not shapes? Why exactly this many?
\solref{lat3}
\end{exercise}
\begin{exercise}\label{ex:dual}
Apply each relabelling to $x\op(x\cop y)=x$: the swap, the opposite of $\op$, the opposite of $\cop$, and both
opposites. \solref{dual}
\end{exercise}
}

% Section 6: Synergy: when the whole does more than the parts
\expandafter\newcommand\csname exercises@6\endcsname{%
\exgroup{6}{Synergy: when the whole does more than the parts}%
\begin{exercise}\label{ex:emerg}
Using Exercises~\ref{ex:abs} and~\ref{ex:abs3} (and the mirror image of~\ref{ex:abs3}), show that $x\op x=x$ is
emergent in the absorption pair. \solref{emerg}
\end{exercise}
\begin{exercise}\label{ex:defsyn}
Let $T=\{\,x\cop y=y\op x,\ (x\op y)\op z=x\op(y\op z)\,\}$. Show that $T\models (x\cop y)\cop z=x\cop(y\cop z)$, and that
this is emergent. Why is it nevertheless not an interesting fusion? \solref{defsyn}
\end{exercise}
}

% Section 7: Disguises: synergy that is not a real fusion
\expandafter\newcommand\csname exercises@7\endcsname{%
\exgroup{7}{Disguises: synergy that is not a real fusion}%
\begin{exercise}\label{ex:rb}
A \emph{rectangular band} on $I\times J$ is the operation $(i,j)\op(k,l)=(i,l)$. On $\{0,1\}\times\{0,1\}$, check that it
is associative and idempotent, and that $x\op(y\op x)=x$ and $(x\op y)\op z=x\op z$ hold. \solref{rb}
\end{exercise}
\begin{exercise}\label{ex:31}
With $\op$ the rectangular band of Exercise~\ref{ex:rb} and $x\cop y:=y\op x$, check the four laws
$x\op(x\cop y)=x$, $(x\cop y)\op y=y$, $(x\op x)\op(y\op x)=x$ and $(x\op y)\op(x\op x)=x$. (These four laws are a theory the
search found; see Section~\ref{sec:found}.) \solref{31}
\end{exercise}
\begin{exercise}\label{ex:split}
The theory $\{x\op y=x,\ x\cop y=x\}$ has no bridge, yet it implies the bridge $x\op y=x\cop y$. Why? Which rule
excludes it? \solref{split}
\end{exercise}
\begin{exercise}\label{ex:split2}
Prove the claim of the split-theory example: under $x\op(x\op y)=y$, the law $x\op(y\cop z)=x\op(z\cop y)$ holds if and only
if $y\cop z=z\cop y$. \solref{split2}
\end{exercise}
}

% Section 8: How the search is made
\expandafter\newcommand\csname exercises@8\endcsname{%
\exgroup{8}{How the search is made}%
\begin{exercise}\label{ex:pairs}
How many theories of exactly two laws can be formed from the $71{,}369$ laws? Compare with the $4{,}222$ theories the
project has verified. \solref{pairs}
\end{exercise}
}


\begin{document}

\begin{center}
{\Large\bfseries Two operations and a machine that looks for surprises\par}
\medskip
{\large An introduction from scratch to the \proj{} project\par}
\medskip
22 November 2026
\end{center}

\bigskip
\noindent\textbf{About this text.} It assumes nothing beyond sets and functions. It builds, step by step, the few
ideas needed to understand what the \proj{} project does: binary operations, laws, theories and their models, and what
it means for laws to \emph{imply} other laws. It then explains the project's question. Can a machine, searching
without help, find small sets of laws in two operations whose combination does more than the parts? It also
explains how that search is built, what it has found so far, and why the question could matter.
\either{Most sections end with exercises, and the solutions are at the end.}{Exercises and their solutions are kept
out of the way, in two appendices at the end.} Numbers about the project are as of 22 November 2026; the
repository's \file{REPORT.md} is the authoritative record, with the command behind every number.

\medskip
\renewcommand{\contentsname}{Contents}
\tableofcontents

\partline{I}{From scratch}

\section{Sets and functions}

A \emph{set} is a collection of objects, its \emph{elements}. We write $a\in A$ for ``$a$ is an element of $A$'' and
list small sets in braces: $\{0,1,2\}$. A set is \emph{finite} if it has $n$ elements for some natural number $n$;
the set $\N=\{0,1,2,\dots\}$ of natural numbers is \emph{infinite}.

The \emph{Cartesian product} $A\times B$ is the set of ordered pairs $(a,b)$ with $a\in A$ and $b\in B$. If $A$ has
$m$ elements and $B$ has $n$, then $A\times B$ has $mn$.

A \emph{function} $f\colon A\to B$ assigns to every $a\in A$ exactly one $f(a)\in B$. If $A$ has $m$ elements and $B$
has $n$, there are $n^m$ functions from $A$ to $B$: each of the $m$ elements independently picks one of $n$ values.

\begin{definition}
A function $f\colon A\to B$ is
\emph{injective} if different inputs give different outputs ($f(a)=f(a')$ implies $a=a'$);
\emph{surjective} if every $b\in B$ is an output ($b=f(a)$ for some $a$);
\emph{bijective} if it is both. A bijection has an \emph{inverse} $f^{-1}\colon B\to A$ with $f^{-1}(f(a))=a$ and
$f(f^{-1}(b))=b$.
\end{definition}

\begin{proposition}\label{prop:finite}
If $A$ is finite, every injective function $f\colon A\to A$ is bijective. On an infinite set this fails:
$n\mapsto n+1$ on $\N$ is injective but not surjective, since $0$ is not an output.
\end{proposition}

\begin{proof}
If $A$ has $n$ elements and $f$ is injective, the $n$ outputs $f(a)$ are all different. So there are $n$ different
outputs inside a set of $n$ elements, and every element is an output.
\end{proof}

This small fact comes back later. Some laws say ``a certain map is injective''. In a finite structure that
already makes it bijective, but in an infinite one it does not. So a statement can be true in every \emph{finite}
structure of a kind and still not follow from the laws.

\inlineexercises{1}

\section{Binary operations}

\begin{definition}
A \emph{binary operation} on a set $A$ is a function $A\times A\to A$. We write it in the middle: $x\op y$ instead of
${\op}(x,y)$. A set with one binary operation is called a \emph{magma}.
\end{definition}

On a finite set we can write the operation as a table, its \emph{Cayley table}: the entry in row $x$ and column $y$
is $x\op y$.

\begin{example}\label{ex:tables}
Four operations on $\{0,1,2\}$:
\begin{center}
\begin{tabular}{c|ccc}
$+_3$ & 0 & 1 & 2\\ \hline
0 & 0 & 1 & 2\\ 1 & 1 & 2 & 0\\ 2 & 2 & 0 & 1
\end{tabular}\qquad
\begin{tabular}{c|ccc}
$\max$ & 0 & 1 & 2\\ \hline
0 & 0 & 1 & 2\\ 1 & 1 & 1 & 2\\ 2 & 2 & 2 & 2
\end{tabular}\qquad
\begin{tabular}{c|ccc}
$x\op y=x$ & 0 & 1 & 2\\ \hline
0 & 0 & 0 & 0\\ 1 & 1 & 1 & 1\\ 2 & 2 & 2 & 2
\end{tabular}\qquad
\begin{tabular}{c|ccc}
$x\op y=0$ & 0 & 1 & 2\\ \hline
0 & 0 & 0 & 0\\ 1 & 0 & 0 & 0\\ 2 & 0 & 0 & 0
\end{tabular}
\end{center}
They are: addition modulo 3, the maximum, the \emph{left projection} (it ignores its right argument), and a
\emph{constant} operation.
\end{example}

An operation on $n$ elements is a table with $n^2$ cells, each holding one of $n$ values. So there are
$n^{n^2}$ operations on an $n$-element set. That is $16$ on two elements, $19{,}683$ on three, and about $4.3$ billion on
four. Small sets already hold more operations than anyone could look at one by one.

Projections and constant operations are \emph{degenerate}: they do not really combine their two arguments. We will
need to rule them out later.

\inlineexercises{2}

\section{Laws}

To talk about all elements at once we use \emph{variables} $x,y,z,\dots$ and build \emph{terms} from them with the
operation: $x$, $x\op y$, $(x\op y)\op z$, $x\op(y\op x)$, and so on. A \emph{law} (an identity, an equation) is a
pair of terms written $s=t$, for example
\begin{align*}
(x\op y)\op z &= x\op (y\op z) && \text{(associativity)},\\
x\op y &= y\op x && \text{(commutativity)},\\
x\op x &= x && \text{(idempotence)}.
\end{align*}
A law \emph{holds} in a magma if the two sides are equal for \emph{every} choice of elements for the variables. So
associativity holds in $(A,\op)$ when $(a\op b)\op c=a\op(b\op c)$ for all $a,b,c\in A$.

On a finite set, checking a law is a finite job. A law with $k$ variables on an $n$-element set needs $n^k$ checks:
try every assignment and compare the two sides. A law fails as soon as one assignment gives different sides. That
assignment is a \emph{counterexample}.

The \emph{order} of a law is the number of operation symbols in it. Associativity has order $4$ (two on each side),
commutativity $2$, idempotence $1$. The project works with all laws of order at most $4$.

Two laws that differ only by the names of their variables say the same thing: $x\op y=y\op x$ and $u\op v=v\op u$.
So do $s=t$ and $t=s$. Laws are therefore counted \emph{up to renaming}: each law is written once, with its
variables named $x,y,z,\dots$ in order of first appearance. In one operation, for example, there are exactly five
laws of order~1 up to renaming\either{ (Exercise~\ref{ex:order1})}{: $x\op x=x$, $x\op x=y$, $x\op y=x$, $x\op y=y$
and $x\op y=z$}.

\inlineexercises{3}

\section{Theories, models and implication}

\begin{definition}
A \emph{theory} is a finite set $T$ of laws. A \emph{model} of $T$ is a structure in which every law of $T$ holds.
$T$ \emph{implies} a law $e$, written $T\models e$, if $e$ holds in \emph{every} model of $T$, finite or infinite.
\end{definition}

There are exactly two ways to settle a question ``does $T\models e$?''.
\begin{itemize}
\item \textbf{Refute it} with a \emph{countermodel}: a model of $T$ in which $e$ fails. One is enough.
\item \textbf{Prove it} by equational reasoning. Start from the laws of $T$, substitute terms for variables, and
  replace a part of a term by something equal to it, until you reach $e$.
\end{itemize}

\begin{example}[an equational proof]
Let $T=\{x\op y=x\}$. We show $T\models(x\op y)\op z=x\op(y\op z)$. Substituting the term $x\op y$ for $x$, and $z$ for
$y$, in the law of $T$ gives $(x\op y)\op z=x\op y$. The law itself then gives $x\op y=x$. Substituting $y\op z$ for $y$
gives $x\op(y\op z)=x$. Both sides equal $x$, so they are equal. Each step uses only the law, substitution, and
replacing equals by equals.
\end{example}

\begin{example}[a countermodel]
Does $\{x\op x=x\}$ imply $x\op y=y\op x$? No. The left projection on $\{0,1\}$ is idempotent, and $0\op1=0\neq1=1\op0$.
\end{example}

Birkhoff's completeness theorem (1935) says these are the only two cases: if $T\models e$, then an equational proof
exists. The class of all models of a theory is called a \emph{variety}. Lattices, rings, and groups (with the
identity and the inverse written as operations of their own) are all varieties.

The trouble is that neither side is guaranteed to finish. A proof may be very long. A countermodel may need many
elements, and sometimes every countermodel is infinite. In general no algorithm decides every question of this
kind; closely related problems, like the word problem for semigroups, are undecidable. Computers help on both
sides:
\begin{itemize}
\item \emph{finite-model finders} (Mace4, Vampire's finite model builder) search for countermodels;
\item \emph{theorem provers} (E, Vampire, Prover9, Twee) search for proofs.
\end{itemize}
A prover can also \emph{saturate}: it runs out of new consequences without reaching $e$. A saturation proves that $e$
does \emph{not} follow, even when every countermodel is infinite.

The Equational Theories Project (ETP, 2024--2025) did this at scale for one operation. It decided all of the about
22 million implications between pairs of single laws of order at most $4$ (there are $4{,}694$ such laws up to
renaming), with proofs and countermodels checked in the proof assistant Lean~\cite{etp}.

\inlineexercises{4}

\section{Two operations}

From now on a structure is a set with \emph{two} binary operations, written $x\op y$ and $x\cop y$. Laws may use
both. Many of the most studied structures in algebra are of this kind.
\begin{itemize}
\item \textbf{Lattices.} $\max$ and $\min$ on a chain such as $\{0,1,2\}$, or union and intersection of sets. The
  laws: both operations associative, commutative and idempotent, plus the two \emph{absorption laws}
  $x\op(x\cop y)=x$ and $x\cop(x\op y)=x$.
\item \textbf{Rings.} $+$ and $\times$ on $\Z$, linked by distributivity $x\times(y+z)=(x\times y)+(x\times z)$. That
  law has order $5$, beyond what the project enumerates.
\item \textbf{Quasigroups with their divisions.} A \emph{left quasigroup} has an operation $\op$ in which every
  equation $a\op u=b$ has exactly one solution $u$. Writing that solution as $a\cop b$ gives two laws,
  $x\op(x\cop y)=y$ and $x\cop(x\op y)=y$. Here $\cop$ is just the division of $\op$, written as a second
  operation.
\end{itemize}

\begin{definition}
A law that mentions both operations is a \emph{bridge}. The absorption laws are bridges; associativity of $\op$ is
not.
\end{definition}

Bridges are where two theories, one for each operation, get fused into one structure.

\paragraph{Symmetries.} Renaming changes nothing essential. We may \emph{swap} the two operations, or replace an
operation by its \emph{opposite} ($x\op' y:=y\op x$). These three moves (swap, opposite of $\op$, opposite of $\cop$)
generate 8 relabellings. A theory that is mapped to itself by one of them has a \emph{duality}. Lattices are the
classic case: swapping $\max$ and $\min$ gives a lattice again.

\inlineexercises{5}

\partline{II}{The project}

\section{Synergy: when the whole does more than the parts}\label{sec:synergy}

\either{Exercises~\ref{ex:abs} and~\ref{ex:abs3} contain the central idea. Neither absorption law alone forces
$x\op x=x$, but the two together do.}{The absorption laws of the previous section contain the central idea. Together
they force $x\op x=x$. The second law with $y=x$ gives $x\cop(x\op x)=x$. The first law with $y=x\op x$ gives
$x\op\bigl(x\cop(x\op x)\bigr)=x$, and replacing the inner term by $x$ turns this into $x\op x=x$. Neither law alone
does it. On $\{0,1\}$, let $x\cop y=0$ for all $x,y$, and let $\op$ have $0\op0=0\op1=0$, $1\op0=1$ and $1\op1=0$. The first
law holds, but $1\op1\neq1$. Exchanging the roles of the two operations shows the same for the second law.}

\begin{definition}[\cite{paper}]
Let $T$ be a theory. A law $e$ is \emph{emergent} in $T$ if $T\models e$ but no single law of $T$ implies $e$. A
theory with emergent laws has \emph{synergy}.
\end{definition}

Synergy is what makes a combination of laws more than a list. The absorption pair is the reference case: two
bridges, and from them a law (idempotence) that neither one gives. It also shows that synergy can be deep. Ralph
McKenzie found in 1970 a basis of just four laws for lattice theory~\cite{mckenzie1970}. Commutativity and
associativity, which the usual definition lists as axioms, are consequences of those four laws together; no single
one of them implies them.

\emph{The question of the project:} can an automatic search over theories, scored by their emergent laws and checked
by provers, find structure that a mathematician would call interesting? In particular, can it find theories that are
\emph{not already named}? No language model is used anywhere in the search or its scoring, so what it finds is
found, not recalled.

\inlineexercises{6}

\section{Disguises: synergy that is not a real fusion}

A search that rewards emergent laws will find every cheap way to produce them. The project learned this the hard
way. Four kinds of theory look like fusions and are not.
\begin{itemize}
\item \textbf{Collapse.} If a theory forces an operation to be a projection or a constant, thousands of laws become
  true at once. The project requires a model in which both operations are essentially binary and different.
\item \textbf{Definitions.} If one law says $x\cop y=y\op x$, or any other way of writing $\cop$ as a term in $\op$,
  then $\cop$ is not a second operation at all\exnote{ (Exercise~\ref{ex:defsyn})}. Such theories are filtered out.
\item \textbf{Split theories.} Some theories are really a theory of $\op$ plus a separate theory of $\cop$; their
  bridges only look like bridges. For example, if a law makes every map $y\mapsto x\op y$ a bijection, an
  apparent bridge can turn out to be a law about one operation alone.
\item \textbf{Divisions in disguise.} The quasigroup laws of the previous section, possibly relabelled, make $\cop$
  the division of $\op$. The result is a one-operation structure written with two symbols.
\end{itemize}
The filters for the last two were added after the search kept ranking such theories at the top.

\begin{example}[a split theory]
Let $T=\{\,x\op(x\op y)=y,\ x\op(y\cop z)=x\op(z\cop y)\,\}$. The second law is a bridge. But the first law says that the
map $y\mapsto x\op y$ undoes itself, so it is injective. Then the second law can be ``cancelled'': it says exactly
$y\cop z=z\cop y$. So $T$ has the same models as $\{x\op(x\op y)=y,\ y\cop z=z\cop y\}$, a law about $\op$ alone next to a law
about $\cop$ alone. The bridge was a disguise\exnote{ (Exercise~\ref{ex:split2})}.
\end{example}

\inlineexercises{7}

\section{How the search is made}

\paragraph{The laws.} A program lists every law of order at most $4$ in two operations, up to renaming the
variables and swapping the two sides: $71{,}369$ laws, of which $61{,}983$ are bridges. For one operation the same
program gives $4{,}694$ laws, the ETP's count, as a check.

\paragraph{A cheap guess at implication.} Proving $T\models e$ is expensive, so the search first guesses. It keeps a
fixed collection of $12{,}509$ small structures on 2 to 6 elements: all pairs of operations on 2 elements, affine
operations modulo $n$, lattices, group-derived operations, random tables. It records, once, which of the $71{,}369$
laws hold in each of them. To guess what $T$ implies, it looks at the structures that satisfy $T$ and collects the
laws true in all of them. \either{Exercise~\ref{ex:corpus} shows why this is only a guess: a law can hold in every
structure you happen to have and still not follow.}{This is only a guess: a law can hold in every structure you
happen to have and still not follow. For example, $\max$ and the left projection on $\{0,1\}$ are both idempotent and
both associative. Yet $x\op y=2x+2y\bmod 3$ on $\{0,1,2\}$ is idempotent and not associative.} Every countermodel the verifier finds later is added to the collection,
so the guesses improve.

\paragraph{The search.} The search is a \emph{quality-diversity} algorithm, MAP-Elites~\cite{mapelites}. Theories are
sorted into cells by a few features (number of laws, number of dualities, number of bridges, how many models, the
order of the gained laws), and each cell keeps the best theory found so far. New theories come from old ones: replace
a law, add or drop one, combine two theories. Keeping the best of every cell, not just the best overall, keeps the
search from collapsing onto one idea. One variant rewards theories with a duality (the \emph{duality prior}).

\paragraph{Verification.} Every theory that wins a cell goes to a queue of verifiers. They work in stages.
\begin{enumerate}
\item Search for small models (2 to 5 elements), check that the laws are independent, and look for countermodels
  (up to 6 elements) to each guessed emergent law.
\item Run the prover E with a time limit (and the solver Z3). A proof confirms the law. A saturation refutes it, and
  is then rechecked by an independent program. Vampire's finite model builder tries the claims E leaves open.
\item For claims still open: longer runs, and finite-model search with Mace4 up to 10 elements.
\end{enumerate}
Only claims with a proof count. The smallest models of the top candidates are certified once more in Lean, a proof
assistant, by brute-force evaluation.

\paragraph{A worked example.} Follow the absorption pair $T=\{x\op(x\cop y)=x,\ x\cop(x\op y)=x\}$ through the
pipeline.
\begin{enumerate}
\item \emph{Guess.} Among the 12,509 structures, those satisfying $T$ include the lattices and many projection
  pairs. The laws true in all of them, minus the laws guessed for either law of $T$ alone, are the guessed emergent
  laws: $x\op x=x$, $x\cop x=x$, $x\op x=x\cop x$, and others.
\item \emph{Models.} The search for small models finds the 2-element lattice\exnote{ (Exercise~\ref{ex:abs2})}, so
  the theory has a non-degenerate model.
\item \emph{Refute or prove.} No countermodel to $x\op x=x$ exists, and E finds the two-step proof of
  \either{Exercise~\ref{ex:abs}}{Section~\ref{sec:synergy}} at once. The claim is confirmed.
\item \emph{Score.} The theory scores for each proved emergent law, weighted by its difficulty and by the richness of
  its small models.
\item \emph{Identify.} The post hoc checks note that every lattice satisfies $T$: $T$ is a \emph{fragment} of lattice
  theory. It is a rediscovery, and it would be counted as one.
\end{enumerate}

\paragraph{After the search.} The best theories become \emph{candidates}. Before a candidate is looked at, it is
checked against known structures under all 8 relabellings, tested for splits and hidden definitions, and searched
for in the literature. A candidate is not a result. ``New'' needs a literature check, an equivalence check, and
ideally someone who knows the field.

\inlineexercises{8}

\section{What has been found so far}\label{sec:found}

As of 22 November 2026, honestly stated:
\begin{itemize}
\item \textbf{Rediscovery.} The search keeps finding lattice theory in pieces. It rediscovered two of McKenzie's four
  laws, twelve different bases for lattices, and Fried's weakly associative lattices. Against a list of 13 known
  two-operation structures fixed in advance, it found 7. With the duality prior, a third of the verified theories
  are fragments of a named theory, against 8\% without it.
\item \textbf{Prior art from a distant field.} The core of one top candidate is $(x\op y)\cop(y\op x)=x$ and
  $(x\cop y)\op(y\cop x)=x$. This turned out to be known from the study of the Yang--Baxter equation in physics and
  combinatorics: \emph{$\Delta$-bijective} algebras~\cite{bonatto}. The search was not aimed at that field.
\item \textbf{A non-associative relative of a known structure.} The search found the four laws $x\op(x\cop y)=x$,
  $(x\cop y)\op y=y$, $(x\op x)\op(y\op x)=x$ and $(x\op y)\op(x\op x)=x$\exnote{ (Exercise~\ref{ex:31})}. Add associativity and
  you get exactly the rectangular skew lattices of Leech~\cite{leech2020}, up to notation. Without associativity the
  theory has more models, and no source for it was found. It remains a candidate.
\item \textbf{A diagonal attractor.} 97\% of the scored theories prove an emergent law in one variable, such as
  $x\op x=x$ or $x\op x=x\cop x$, and none proves emergent associativity. The search had learned to fuse cheaply
  through the diagonal. A new run (X4) counts only laws that do not come from the diagonal.
\item \textbf{Scale.} All runs together verified $4{,}222$ theories. Sets of up to six of the $71{,}369$ laws number
  about $1.8\times10^{26}$. The search has seen a speck of the space, and mostly the corner humans know best.
\item \textbf{Open problems.} Four implications were shown false although every countermodel is infinite (a
  prover's saturation settles them). One question from the early runs is still open: whether four particular laws
  imply $x\op x=x\cop x$ (``S12'' in the README).
\end{itemize}
No theory found so far is claimed to be new.

\section{Why it could matter}

\begin{itemize}
\item \textbf{Discovery that can be checked completely.} In equational logic every claim ends in a proof or a model,
  and both can be checked by a machine. That makes it one of the few places where automated mathematical discovery
  can be measured without argument.
\item \textbf{Conjunctions are unmapped.} The ETP mapped implications between \emph{single} laws of one operation.
  What sets of laws do together, especially in two operations, is a far larger territory, and most of it is
  unexplored.
\item \textbf{Fusion is how structures arise.} Lattices, rings and quasigroups are all fusions of two operations
  through a few bridges. A search for synergy is a search for the next such structure, if there is one.
\item \textbf{New, not just next.} The hope is a theory like a new genre of music: two existing traditions joined by
  a bridge into something with features neither had, not reachable by small steps from what is known. Such a theory
  would work (have proofs and many models) without being clearly connected to anything named. Whether the search
  can find one is the open question.
\item \textbf{By-products.} Non-implications whose countermodels are all infinite are hard test problems for
  provers, and the rediscovery measurements calibrate how good automatic search is at algebra.
\end{itemize}

\section{Further reading}

\begin{itemize}
\item S.~Burris and H.~P.~Sankappanavar, \emph{A Course in Universal Algebra} (Springer, 1981; free online edition):
  varieties, Birkhoff's theorems.
\item The Equational Theories Project~\cite{etp}: the one-operation map, with its data and Lean proofs.
\item W.~McCune and R.~Padmanabhan, \emph{Automated Deduction in Equational Logic and Cubic Curves} (LNCS 1095,
  1996): the closest earlier method (enumerate laws, filter by small structures, prove).
\item The project's own records: the README; the report, with every number and the command behind it
  (\file{REPORT.md}); the preregistrations (\file{PREREGISTRATION.md}); and the draft paper in \file{tex/}.
\end{itemize}

\appendix
\ifreading
\clearpage
\section*{Exercises}
\addcontentsline{toc}{section}{Exercises}

Exercises for each section of the text, in order. The solutions follow in the next appendix.
\csname exercises@1\endcsname
\csname exercises@2\endcsname
\csname exercises@3\endcsname
\csname exercises@4\endcsname
\csname exercises@5\endcsname
\csname exercises@6\endcsname
\csname exercises@7\endcsname
\csname exercises@8\endcsname
\clearpage
\fi
\section*{Solutions}
\addcontentsline{toc}{section}{Solutions}

The numeric answers are computed by \file{intro/checks.py} (output in \file{intro/checks.out}).

\solution{prod}
$\{(0,0),(0,1),(0,2),(1,0),(1,1),(1,2)\}$: $2\cdot 3=6$ elements.

\solution{funcs}
$2^3=8$ functions. None is injective: three inputs cannot have three different outputs in a two-element set. Six
are surjective: all except the two constant functions.

\solution{inj}
There are $3!=6$ injective functions $\{0,1,2\}\to\{0,1,2\}$ (the permutations), and all six are surjective, as the
proposition says.

\solution{count}
A table on $n$ elements has $n^2$ cells with $n$ choices each: $2^4=16$ and $3^9=19{,}683$.

\solution{idem2}
Idempotence fixes the diagonal ($0\op0=0$, $1\op1=1$) and leaves $0\op1$ and $1\op0$ free: $2^2=4$ operations. They
are $\min$, $\max$, the left projection and the right projection.

\solution{max}
Associativity, commutativity, idempotence and $x\op(x\op y)=x\op y$ hold; $x\op y=x$ fails (for example
$\max(0,1)=1\neq0$). The element $0$ is an identity: $\max(0,x)=x$.

\solution{lp}
$(x\op y)\op z=x\op z=x$ and $x\op(y\op z)=x$, so associativity holds; $x\op x=x$; but $0\op1=0$ while $1\op0=1$, so
commutativity fails.

\solution{cost}
Three variables: $3^3=27$ assignments on 3 elements and $6^3=216$ on 6.

\solution{order1}
Write the product on the left. It is $x\op x$ or $x\op y$. The variable on the other side is either one already used or
a new one: $x\op x=x$, $x\op x=y$, $x\op y=x$, $x\op y=y$, $x\op y=z$. That makes five, and the project's law table
lists the same five. The first is idempotence. The others are degenerate: $x\op y=x$ is the left projection, and
$x\op x=y$ and $x\op y=z$ only hold on a one-element set.

\solution{abs}
The second law with $y=x$ gives $x\cop(x\op x)=x$. The first law with $y=x\op x$ gives $x\op\bigl(x\cop(x\op x)\bigr)=x$.
Replace the inner term using the first step: $x\op x=x$.

\solution{abs2}
Nine pairs. Writing L and R for the left and right projections, $(\op,\cop)$ is one of $(\min,\max)$, $(\max,\min)$,
$(\min,\mathrm{L})$, $(\max,\mathrm{L})$, $(\mathrm{L},\min)$, $(\mathrm{L},\max)$, $(\mathrm{L},\mathrm{L})$,
$(\mathrm{L},\mathrm{R})$, $(\mathrm{R},\mathrm{L})$. Only the first two are non-degenerate: the 2-element lattice, in
its two labellings (exchanging $0$ and $1$ turns one into the other). The other seven use a projection. This is why the project demands a non-degenerate model.

\solution{abs3}
Take $x\cop y=0$ for all $x,y$, and let $\op$ have $0\op0=0$, $0\op1=0$, $1\op0=1$, $1\op1=0$. Then $x\op(x\cop y)=x\op0=x$
for both values of $x$, but $1\op1=0\neq1$.

\solution{corpus}
No. Your two structures agree with ``idempotent implies associative'' only by accident. The operation
$x\op y=2x+2y\bmod 3$ is idempotent ($4x\equiv x$) but not associative: $(0\op0)\op1=0\op1=2$ while
$0\op(0\op1)=0\op2=1$. No counterexample exists on 2 elements; the smallest has 3. This is exactly the situation
of the search's cheap guesses.

\solution{lat3}
Six. A lattice on three elements must be a chain, and each of the $3!=6$ orders of $\{0,1,2\}$ as a chain gives
one pair $(\text{join},\text{meet})$. The swapped pair is the reverse chain, which is already among the six. (By
brute force: 9 of the $19{,}683$ operations are semilattices, and 6 pairs of them satisfy the absorption laws.)

\solution{dual}
Swap: $x\cop(x\op y)=x$. Opposite of $\op$: $(x\cop y)\op x=x$. Opposite of $\cop$: $x\op(y\cop x)=x$. Both
opposites: $(y\cop x)\op x=x$. (Checked by brute force: the law holds in a structure exactly when each dual form
holds in the relabelled structure, for all 256 pairs on $\{0,1\}$.)

\solution{emerg}
$T\models x\op x=x$ by Exercise~\ref{ex:abs}. The model of Exercise~\ref{ex:abs3} satisfies the first law and not
$x\op x=x$, so the first law alone does not imply it. For the second law alone, swap the roles: let $\op$ be the
constant $0$ and $\cop$ the table of Exercise~\ref{ex:abs3}. Then $x\cop(x\op y)=x\cop 0=x$ holds, but
$1\op1=0\neq1$.

\solution{defsyn}
$(x\cop y)\cop z=z\op(y\op x)$ and $x\cop(y\cop z)=(z\op y)\op x$, which are equal by associativity of $\op$. Neither law
alone suffices: the definition alone allows a non-associative $\op$, and associativity alone says nothing about
$\cop$. (By brute force on 3 elements, every associative $\op$ has an associative transpose.) It is not interesting
because $\cop$ is defined by $\op$: the ``theory of two operations'' is a theory of one operation, written twice.

\solution{rb}
With elements coded as pairs, $\bigl((i,j)\op(k,l)\bigr)\op(m,n)=(i,n)=(i,j)\op\bigl((k,l)\op(m,n)\bigr)$, and
$(i,j)\op(i,j)=(i,j)$. Then $x\op(y\op x)$ has the row of $x$ and the column of $x$, so it equals $x$. And
$(x\op y)\op z$ has the row of $x$ and the column of $z$, as does $x\op z$.

\solution{31}
With $x=(i,j)$, $y=(k,l)$: $x\cop y=y\op x=(k,j)$. Then $x\op(x\cop y)=(i,j)\op(k,j)=(i,j)=x$, and
$(x\cop y)\op y=(k,j)\op(k,l)=(k,l)=y$. The last two laws hold because the band is idempotent and
$x\op(y\op x)=x=(x\op y)\op x$. So every rectangular band with its reverse is a model. The search's theory has other,
non-associative models as well.

\solution{split}
Both laws force projections, so both operations are the left projection and are equal: $x\op y=x=x\cop y$. The
bridge is a by-product of collapse. The rule ``require a model with two different, essentially binary
operations'' excludes the theory.

\solution{split2}
If $y\cop z=z\cop y$, the two sides of the bridge are equal. Conversely, apply the map $u\mapsto x\op u$ to both sides of
$x\op(y\cop z)=x\op(z\cop y)$. By $x\op(x\op u)=u$ this gives back $y\cop z=z\cop y$. (By brute force, the two theories have
the same models on 2 and on 3 elements: 32 and 46,656.)

\solution{pairs}
$\binom{71369}{2}=2{,}546{,}731{,}396$, about $2.5\times10^9$ pairs, and about $1.1\times10^{18}$ sets of four. The
$4{,}222$ verified theories are less than two in a million of the pairs alone.

\begin{thebibliography}{9}
\bibitem{paper} The \proj{} project, draft paper \file{tex/paper.pdf} and \file{REPORT.md} in the repository.
\bibitem{etp} M.~Bolan, J.~Breitner, J.~Brox, N.~Carlini et al., \emph{The Equational Theories Project: Advancing
  Collaborative Mathematical Research at Scale}, arXiv:2512.07087 (2025).
\bibitem{mckenzie1970} R.~McKenzie, Equational bases for lattice theories, \emph{Math.\ Scand.}\ 27 (1970) 24--38.
\bibitem{mapelites} J.-B.~Mouret and J.~Clune, Illuminating search spaces by mapping elites, arXiv:1504.04909 (2015).
\bibitem{bonatto} M.~Bonatto, M.~Kinyon, D.~Stanovsk\'y and P.~Vojt\v{e}chovsk\'y, Involutive latin solutions of the
  Yang--Baxter equation, arXiv:1910.02148 (2019).
\bibitem{leech2020} J.~Leech, \emph{Noncommutative Lattices: Skew Lattices, Skew Boolean Algebras and Beyond},
  University of Primorska Press (2020).
\end{thebibliography}

\end{document}
